\subsection{图的庞加莱群}

设 G 为图。设 a 为 G 的顶点；群胚 \(\mathrm{Grp}(G)\) 在 a 处的迷向群称为 G 在 a 处的庞加莱群，记作 \(\pi_{1}(G,a)\) 。设 c 为 G 中的道路类，a 为其起点，b 为其靶。映射 \(\operatorname{Int}(c): \pi_{1}(G,b)\to\pi_{1}(G,a)\) （由 \(c'\mapsto cc'c^{-1}\) 定义）是群同构。设 \(\varphi:G\to H\) 为图态射。以 \(\theta_{G}:G\to\mathrm{Grp}(G)\) 与 \(\theta_{H}:H\to\mathrm{Grp}(H)\) 表示典范态射。若以 \(\overline{\varphi}\) 表示唯一的群胚态射 \(\mathrm{Grp}(G) \to \mathrm{Grp}(H)\) （它满足 \(\overline{\varphi}\circ\theta_{G}=\theta_{H}\circ\varphi\) ），则群同态 \(\overline{\varphi}_{a}:\pi_{1}(G,a)\to\pi_{1}(H,\varphi(a))\) 记作 \(\pi_{1}(\varphi,a)\) 。

设 G 为连通图，S 为 G 的一个定向，A 为 G 的极大定向树（II, 第 157 页, 命题 1）。给定 G 的顶点 \(a\) 与 \(b\) ，存在唯一的道路类 \(\gamma_{a,b}\) （在与 A 相伴的图 \(\widetilde{\mathrm{A}}\) 中），它连接 \(a\) 与 \(b\) （II, 第 174 页, 命题 5 的推论 4）。若 \(a\) 、\(b\) 与 \(c\) 是 G 的顶点，则有 \(\gamma_{a,b}\gamma_{b,c} = \gamma_{a,c}\) ，这两个道路类都等于连接 \(a\) 与 \(c\) 的唯一道路类（在 \(\widetilde{\mathrm{A}}\) 中）。

命题 9. —— 设 a 为 G 的顶点。存在唯一的同态 \(\lambda\) ，它从自由群 \(\mathrm{F}(\mathrm{S}-\mathrm{Fl}(\mathrm{A}))\) 到 \(\pi_1(G, a)\) ，使得

\[
\lambda (f) = \gamma_ {a, o (f)} \cdot f \cdot \gamma_ {t (f), a}\tag{1}
\]

对 \(f \in \mathrm{S} - \mathrm{Fl}(\mathrm{A})\) 成立。同态 \(\lambda\) 是群同构。

以 L 表示群 \(\mathrm{F}(\mathrm{S}-\mathrm{Fl}(\mathrm{A}))\) 。满足关系式 (1) 的同态 \(\lambda\colon L\to\pi_{1}(\mathrm{G},a)\) 的存在性与唯一性由 A, I, 第 85 页, 命题 8 得出。设 L 为与群 L 相伴的群胚（II, 第 163 页, 例 1）；以 s 表示其唯一的顶点。存在唯一的群胚态射 \(\mu\colon\mathrm{Grp}(\mathrm{G})\to\mathcal{L}\) ，使得 \(\mu(\theta_{\mathrm{G}}(f))=f\) （对每条箭 \(f\in\mathrm{S}-\mathrm{Fl}(\mathrm{A})\) ）且 \(\mu(\theta_{\mathrm{G}}(f))=e_{s}\) （对每条箭 \(f\in\mathrm{Fl}(\mathrm{A})\) ）（II, 第 173 页, 命题 5）。对 \(f\in\mathrm{S}-\mathrm{Fl}(\mathrm{A})\) ，有 \(\mu(\theta_{\mathrm{G}}(f))=f\) ；因此同态 \(\mu_{a}\circ\lambda\) 是群 L 的恒同同构（A, I, 第 85 页, 命题 8）。由此推出 \(\lambda\) 是单射的。

设 \(\varpi_{\mathrm{G}} / \mathrm{A}\) 为由 \(\varpi_{\mathrm{G}}\) 经缩并 A 的箭而得的群胚，设 \(p\colon \varpi_{\mathrm{G}}\to \varpi_{\mathrm{G}} / \mathrm{A}\) 为典范群胚态射。它是满射的，且群同态 \(p_a\) （由 \(p\) 通过取迷向群而得）是同构（II, 第 178 页, 命题 8 的推论 2）。由于假定图 G 连通且 A 是它的极大树，群胚 \(\varpi_{\mathrm{G}} / \mathrm{A}\) 有唯一的顶点（II, 第 176 页）。此外，\(\varpi_{\mathrm{G}}\) 由 \(\theta_{\mathrm{G}}(\mathrm{G})\) 生成；因此 \(\varpi_{\mathrm{G}} / \mathrm{A}\) 由诸圈 \(p(f)\) （对 \(f\in \mathrm{S}\) ）生成，且群 \((\varpi_{\mathrm{G}} / \mathrm{A})_{p(a)}\) 由形如 \(p(\theta_{\mathrm{G}}(f))\) 的元素（对 \(f\in \mathrm{S} - \mathrm{Fl}(\mathrm{A})\) ）生成。因此同态 \(p_a\circ \lambda\) 是满射的。从而 \(\lambda\) 是满射的。

注. —— 1) 同态 \(\mu_{a}\) 是 \(\lambda\) 的逆同构。

\begin{enumerate}
\def\labelenumi{\arabic{enumi})}
\setcounter{enumi}{1}
\tightlist
\item
  存在唯一的同态 \(\lambda\colon\mathrm{F}(\mathrm{Fl}(\mathrm{G}))\to\pi_{1}(\mathrm{G},a)\) ，它对一切 \(f\in\mathrm{Fl}(\mathrm{G})\) 由关系式 (1) 定义。由命题 8 推出，同态 \(\lambda\) 是满射的，其核是 \(\mathrm{F}(\mathrm{Fl}(\mathrm{G}))\) 的最小正规子群，它包含元素 \(f\) （对 \(f\in\mathrm{Fl}(\mathrm{A})\) ）以及元素 \(f\cdot\overline{f}\) （对 \(f\in\mathrm{Fl}(\mathrm{G})\) ）。
\end{enumerate}

